\subsection{Изменение кинетической энергии}
%%%%%%%%%%%%%%%%%%%%%%%%%%%%%%%%%%%%%%%%%%%%%%%%%%%%
%%%%%%%%%%%%%%%%%%%%%%%%%%%%%%%%%%%%%%%%%%%%%%%%%%%%
 \z{Частица массы $m$ со скоростью $v$ влетает в область действия тормозящей силы $F$ под углом $\alpha$ к направлению этой силы. Под каким углом $\beta$ к направлению силы $F$ она вылетит из этой области? Ширина области действия силы $l$. При каком условии частица не сможет пересечь эту область?}
%
{$\sin\beta = \frac{\sin\alpha}{\sqrt{1 - \frac{2Fl}{mv^2}}} \quad \text{при} \quad Fl > \frac{mv^2}{2}\cos^2\alpha$}
%%%%%%%%%%%%%%%%%%%%%%%%%%%%%%%%%%%%%%%%%%%%%%%%%%%%
\z[{\begin{wrapfigure}[5]{r}{0.3\textwidth}%
		\vspace{-12pt}%
		\begin{flushright}%
 			\includegraphics[width=0.3\textwidth]{7.2.2}%
		\end{flushright}%
	\end{wrapfigure}%
	}]%
	{Груз массой $m$ толкнули вверх по гладкой доске массой $M$ и длиной $l$, шарнирно закреплённой в точке $O$. Доска с горизонтом составляет угол $\alpha = 45\degree$, расстояние $OA = h < l/2$. Какую скорость $v$ нужно сообщить грузу, чтобы нижний конец доски оторвался от пола?}
%
{$v \geq \sqrt{2g\left(l - h + \frac{M(l - 2h)}{2m}\right)\sin\alpha}, \quad \text{если} \quad \frac{M}{m} < \frac{2h}{l - 2h}, \quad \text{иначе решений нет}$}
%%%%%%%%%%%%%%%%%%%%%%%%%%%%%%%%%%%%%%%%%%%%%%%%%%%% 
 \z{Какова работа силы трения за один оборот аэросаней, движущихся по вертикальной круговой дорожке? Скорость саней постоянна и равна $v$, масса саней $m$, коэффициент трения $\mu = 0.2$.}
%
{$A = 2\pi\mu mv^2$}
%%%%%%%%%%%%%%%%%%%%%%%%%%%%%%%%%%%%%%%%%%%%%%%%%%%% 
 \z{По плоскости с углом наклона $\alpha = 45\degree$ соскальзывает шайба и в конце спуска ударяется о стенку, перпендикулярную наклонной плоскости. На какую высоту $h$ снова поднимется шайба по плоскости, если первоначально она находилась на высоте $H = 0.6$ м? Коэффициент трения $\mu = 0.2$.}
%
{$h = H \cdot \frac{1 - \mu\cot\alpha}{1 + \mu\cot\alpha}$}
%%%%%%%%%%%%%%%%%%%%%%%%%%%%%%%%%%%%%%%%%%%%%%%%%%%%
%%%%%%%%%%%%%%%%%%%%%%%%%%%%%%%%%%%%%%%%%%%%%%%%%%%%	
 \z{Брусок, двигавшийся по горизонтальной поверхности стола со скоростью $v_0$, сталкивается с неподвижным бруском вдвое большей массы. На какое расстояние $d$ разъедутся бруски после столкновения? Удар упругий, центральный. Коэффициенты трения брусков о стол одинаковы и равны $\mu = 0.1$.}
 %
{$L = \frac{5v_0^2}{18\mu g}$}